\section{讨论}

TE-KG 处理了人类 TE 研究中的一个实际整合问题：同一命名元件的证据分散在家族参考资源、基因组注释、表达矩阵和生物医学论文中。该资源通过共享的 TE 身份和相互连接的访问流程连接这些证据层。其主要贡献并不是单独使用图技术，而是将带 PMID 的文献图谱与分类、代表性序列和基因组记录、三类表达情境及特定情境下的共表达视图结合，同时使每个证据层的解释边界保持可见。

这种定位有意与专业资源形成互补。Repbase 和 Dfam 提供更深入的重复序列家族参考与模型 \citep{Bao2015Repbase,Wheeler2013Dfam}。dbRIP 和 euL1db 以 TE-KG 目前无法提供的分辨率描述插入多态性和样本特异性 L1HS 记录 \citep{Wang2006dbRIP,Mir2015euL1db}。HervD Atlas 对 HERV-疾病关联进行更深入的人工整理，TE-SCALE 则以超过本文批量数据集的规模提供癌症单细胞表达和共表达信息 \citep{Li2024HervDAtlas,Meng2026TESCALE}。因此，当问题跨越多种证据类型时，TE-KG 最具价值；如果用户需要某一专业领域中分辨率最高的记录，则应优先使用相应专业资源。

文献图谱具有明确的来源追踪优势，也具有明确限制。保留谓词和 PMID 使图中的边可被检查，并允许用户返回原始论文。然而，当前语料库经过自动筛选、受约束的语言模型提取、标准化和人工整理等步骤生成。检索、缩写解析、实体标准化或原文解释均可能产生错误。PMID 的存在说明关系可以追溯，并不证明关系一定正确，也不证明因果性。在提取流程能够支持定量准确率结论之前，还需要完成可复现的人工审核并保存样本标识符。

表达和共表达也需要类似的谨慎。TE 表达估计对重复序列特征和定量方法较为敏感 \citep{Lanciano2020TEExpression}。TE-KG 的三类情境来自不同研究，不能被解释为配对的正常与癌症实验。正常组织元数据还包含重复标识符和五个无法与矩阵匹配的运行记录。共表达分析进一步将数据简化为达到阈值的两两相关和模块。FDR 控制限制保留统计检验中预期的假发现比例，但不能建立调控关系。展示目录又增加了一层筛选，因为它提供的是受限且经批准的子图，而不是全部离线边。

智能界面降低了跨多个存储检索证据的操作成本，但其地位应低于能够直接看到来源的数据库访问。会话范围内的追问支持简短的研究对话，而无需建立持久用户画像。与此同时，回答质量取决于实体解析、插件选择、本地证据可用性和语言模型综合能力。即使底层资源包含有用记录，界面仍可能遗漏相关文献或给出不完整回答。因此，保持 PubMed 链接可见，并使用面向读者的语言表达不确定性，比展示内部推理步骤更重要。

投稿和公开发布前仍有若干要求需要满足。必须记录准确的 RepeatMasker 和 RepBase 版本、获取日期及再分发条款；必须冻结表达数据的“登录号到样本”清单，包括 SRP013565 子集和五条正常组织元数据不匹配记录；还需要稳定的公开网址、版本化代码和数据存档、许可证及更新策略。当前 \emph{Database} 投稿要求资源免费、无需登录，并至少保持两年可用，因此只有本地部署和 Download 页面并不充分。

下一步最有信息量的评估不是增加更多界面截图，而是建立可复现的跨证据层案例。锁定的 L1HS 示例可以检验用户能否从身份信息进入分类、代表性序列或位置、表达、共表达和文献证据，同时不丢失来源信息或夸大结论。再增加一个 HERV 或 DNA 转座子案例，有助于判断该流程能否推广到代表性最强的 LINE 示例之外。这些案例应把证据缺失作为结果的一部分报告，而不是使用未经核实的生成文字进行补充。
